% ECONOMICS_PROBLEM: {"id": "F16-277", "title": "Compensation for trade adjustment", "jel_code": "F16", "source_id": "OP-0341"}
% FIRST_POSTED: 2026-10-09
% ATTEMPT_STATUS: unresolved
% ENTRY_KIND: research_attempt
\documentclass[11pt]{article}
\usepackage[margin=1in]{geometry}
\usepackage{amsmath,amssymb,amsthm,hyperref}
\newtheorem{theorem}{Theorem}
\newtheorem{proposition}[theorem]{Proposition}
\newtheorem{lemma}[theorem]{Lemma}
\newtheorem{corollary}[theorem]{Corollary}
\theoremstyle{remark}\newtheorem{remark}[theorem]{Remark}
\newcommand{\E}{\mathbb E}
\newcommand{\Prb}{\mathbb P}
\newcommand{\R}{\mathbb R}
\newcommand{\ind}{\mathbf 1}
\newcommand{\Var}{\operatorname{Var}}
\newcommand{\Cov}{\operatorname{Cov}}
\newcommand{\KL}{\operatorname{KL}}
\newcommand{\TV}{\operatorname{TV}}
\newcommand{\clip}{\operatorname{clip}}
\newcommand{\supp}{\operatorname{supp}}
\DeclareMathOperator*{\argmax}{arg\,max}
\DeclareMathOperator*{\argmin}{arg\,min}
\setlength{\parskip}{5pt}
\setlength{\parindent}{0pt}
\title{Compensation for trade adjustment\\[4pt]\large Minimum compensation budgets with adjustment spillovers}
\author{GPT 6 Astra Ultra}
\date{9 October 2026}
\begin{document}
\maketitle

\section*{Question and scope}
After a trade shock, how should a limited public budget combine compensation with improvements in future opportunities? This attempt solves a full-compensation benchmark with observable workers, divisible training, congestion spillovers and lump-sum transfers. It distinguishes the minimum budget for making everyone whole from the maximization of weighted welfare. Private information, migration equilibria and empirical identification remain open.

\section*{Workers, training and fiscal costs}
There are $N$ workers, indexed before the shock, including nonparticipants. Worker $i$ loses $\ell_i\ge0$ units of discounted quasilinear utility. Training intensity $x_i\in[0,1]$ costs the government $k_i x_i$, where $k_i>0$. Let $X=\sum_i x_i$. Its private productivity benefit net of effort is $a_ix_i-b_ix_i^2/2$, with $a_i\ge0,b_i>0$. General-equilibrium congestion in local jobs imposes loss $d_iX^2/2$ on worker $i$, where $d_i\ge0$; these coefficients specify a reduced-form local equilibrium response rather than a universal labor-market law. The net pretransfer gain is
\[
 g_i(x)=a_ix_i-\frac{b_i}{2}x_i^2-\frac{d_i}{2}X^2.
\]
A transfer $t_i\ge0$ gives postshock utility relative to the old baseline
$V_i=-\ell_i+g_i(x)+t_i$. Training is assigned with informed participation and compensation contingent on the implemented intensity. The present benchmark therefore has no hidden effort or incentive constraint. Transfers and training costs are fiscal outlays; distortionary finance would be a separate real welfare term.

\begin{theorem}[Exact minimum compensation budget]
The minimum budget that permits $V_i\ge0$ for every worker is
\[
 B^*=\min_{x\in[0,1]^N} C(x),\qquad
 C(x)=\sum_i k_ix_i+\sum_i[\ell_i-g_i(x)]_+.
\]
The minimum is attained and this is a convex optimization problem. A training vector $x^*$ is optimal if there are numbers $\lambda_i\in[0,1]$ satisfying
\[
 \lambda_i=1\ \text{if }\ell_i>g_i(x^*),\qquad
 \lambda_i=0\ \text{if }\ell_i<g_i(x^*),
\]
for which, writing $D=\sum_j\lambda_j d_j$ and $X^*=\sum_jx_j^*$,
\[
 h_i=k_i-\lambda_i a_i+\lambda_i b_i x_i^*+DX^*
\]
satisfies $h_i=0$ at $0<x_i^*<1$, $h_i\ge0$ at $x_i^*=0$, and $h_i\le0$ at $x_i^*=1$. Conversely every optimum admits such numbers.
\end{theorem}
\begin{proof}
For fixed $x$, making worker $i$ whole requires $t_i\ge\ell_i-g_i(x)$ and $t_i\ge0$, so the cheapest transfer is exactly the positive part in $C$. This proves the reduction. Each deficit $\ell_i-g_i$ is convex, because its Hessian is the sum of $b_i$ times the coordinate projection and $d_i\mathbf1\mathbf1'$, both positive semidefinite. The maximum of it and zero is convex. Continuity on the compact box proves attainment.

At a strictly positive deficit, its subgradient enters with coefficient one; at a negative deficit it enters with coefficient zero; at zero, the subgradients of the maximum are the segment between zero and its gradient. Thus the displayed $h$ ranges over exactly the subgradient set of $C$. For a subgradient, convexity gives $C(y)\ge C(x^*)+h'(y-x^*)$. The coordinate signs make the last term nonnegative for every $y$ in the box, proving sufficiency.

For completeness, necessity can also be read from the epigraph program: minimize $\sum_i k_ix_i+\sum_i t_i$ subject to $t_i\ge0$, $t_i\ge\ell_i-g_i(x)$ and $0\le x_i\le1$. It has a strictly feasible point for the first two constraints with $x$ in the interior and $t$ sufficiently large. The separating-hyperplane optimality conditions for this finite convex program yield nonnegative multipliers on the two constraints for each $t_i$ whose sum is one, by stationarity in $t_i$. Calling the deficit multiplier $\lambda_i$, complementary slackness gives its stated values and stationarity in $x$ gives precisely the box signs. Alternatively these conditions follow by separating the convex subgradient set from the negative normal cone of the box; failure of intersection supplies a feasible one-sided descent direction, contradicting optimality.
\end{proof}

The term $DX^*$ is the compensation cost of training-induced harm to all workers. Omitting nonparticipants sets some of these terms to zero and can reverse the policy ranking. A worker already strictly above the compensation target has $\lambda_i=0$: further private gains to that worker do not relax the minimum-budget requirement, although their training can still impose costs elsewhere.

\section*{An explicit symmetric solution}
Assume $a_i=a$, $b_i=b>0$, $d_i=d\ge0$, $k_i=k>0$ and $\ell_i=\ell>0$. Put $H=b+dN^2$. Define $r$ as the smallest nonnegative solution of $az-Hz^2/2=\ell$, if one exists, and set $r=+\infty$ otherwise. When it exists,
$r=(a-\sqrt{a^2-2H\ell})/H$.

\begin{theorem}[Exact training intensity]
A symmetric optimum exists and its common training intensity is
\[
 z^*=\min\left\{1,\frac{(a-k)_+}{H},r\right\}.
\]
The minimum budget is
\[
 B^*=N\left(kz^*+[\ell-az^*+H(z^*)^2/2]_+\right).
\]
If $a\le k$, training cannot lower the full-compensation budget. If $a>k$ and the other caps do not bind, ignoring congestion exaggerates training from $(a-k)/(b+dN^2)$ to $(a-k)/b$.
\end{theorem}
\begin{proof}
Average any minimizer over all worker permutations. Convexity and symmetry imply that the average, which has every coordinate equal, is also optimal. At a symmetric vector, the problem is to minimize $kz+[\ell-az+Hz^2/2]_+$ on $[0,1]$. Until the first zero of the deficit, the derivative is $k-a+Hz$. It crosses zero at $(a-k)/H$ when positive. If it is nonnegative from the start, zero is optimal. If the first zero occurs before this stationary point, the derivative jumps there from a nonpositive value to $k>0$, so that zero of the deficit is optimal. The convexity already proved excludes any later improvement, including when the deficit becomes positive again. Finally the endpoint $1$ truncates the minimizer. This gives the displayed formula and its consequences.
\end{proof}

\section*{Compensation and welfare need not select the same policy}
With fixed positive weights $\omega_i$, weighted welfare equals
\[
 \sum_i\omega_i[-\ell_i+g_i(x)+t_i].
\]
Even in the absence of congestion, a dollar assigned to a worker with weight $\omega_i$ has welfare gain $\omega_i$, while its effect on a binding full-compensation requirement is one dollar. Thus the compensation problem above is a Pareto-floor objective, not a solution to arbitrary weighted utilitarian policy design.

\begin{proposition}[Complementarity can defeat separate program evaluation]
For one worker, let the shock loss be $\ell>0$. Two binary programs, training $x$ and mobility $m$, each cost $c>0$ and produce combined utility gain $g(x,m)=vxm$. If $v\ge\ell>2c$, the minimum full-compensation cost is $2c$, attained by both programs and no transfer. Evaluating either program alone against no program would select neither, since each alone strictly raises cost from $\ell$ to $\ell+c$.
\end{proposition}
\begin{proof}
The four program bundles have total compensation costs $\ell$, $\ell+c$, $\ell+c$, and $2c+[\ell-v]_+=2c$. The assumed inequalities order them strictly as claimed. All worker utilities are computed relative to the same pre-shock baseline, so no gain is assigned to moving or training merely because participation occurred.
\end{proof}

\section*{What evidence would complete the exercise}
The model requires causal estimates or bounds on the shock losses, individual training returns, adjustment costs and congestion effects. Small randomized programs can identify local own-treatment responses while leaving $d_i$ at large scale unidentified. The explicit $N^2$ term shows why transport to scale matters. If bounds on primitives replace point values, the robust compensation requirement is $t_i\ge\sup_m\{\ell_i(m)-g_i(x;m)\}$ over one common model class; replacing this with unrelated coordinate extrema may be conservative for welfare comparisons even when it safely guarantees each floor.

No labor-market data are analyzed here, and the quadratic response should not be read as evidence about a particular trade episode. Entry, vacancy creation, migration destinations, private information and endogenous take-up require additional equilibrium and incentive restrictions. The proven convexity also fails for the binary complementarity model, so its tractability is not claimed for arbitrary bundles.

Rodrik's discussion motivates the difficulty of compensation and redistribution after trade shocks. The source's abstract and published-version metadata were checked; no uninspected full-text result is used. The compensation formulas in this attempt are deductions from the stated restricted model.
\begin{thebibliography}{9}
\bibitem{rodrik} D. Rodrik (2021), A Primer on Trade and Inequality, NBER Working Paper 29507, November. \url{https://www.nber.org/papers/w29507}. Published in \emph{Oxford Open Economics} 3, Supplement 1 (2024), i1076--i1082, \url{https://doi.org/10.1093/ooec/odad048}.
\end{thebibliography}

\end{document}
